% Part: sets-functions-relations
% Chapter: relations
% Section: special-properties

\documentclass[../../../include/open-logic-section]{subfiles}

\begin{document}

\olfileid{sfr}{rel}{prp}
\olsection{সম্পর্কের বিশেষ ধর্ম}

\begin{intro}
কয়েক ধরনের সম্পর্ক এত ঘন ঘন দেখা যায় যে তাদের বিশেষ নাম দেওয়া হয়েছে। যেমন, $\le$ এবং~$\subseteq$ নিজ নিজ ক্ষেত্রের (ধরা যাক, $\le$-এর ক্ষেত্রে $\Nat$ এবং $\subseteq$-এর ক্ষেত্রে $\Pow{A}$) মধ্যে একই ধরনের সম্পর্ক স্থাপন করে। এই সম্পর্কগুলির মিল ও অমিল ঠিক কোথায়, তা বোঝার জন্য আমরা সম্পর্কের কয়েকটি বিশেষ ধর্ম অনুসারে তাদের শ্রেণিবদ্ধ করি। দেখা যায়, এই ধর্মগুলির কয়েকটির সমন্বয় বিশেষ গুরুত্বপূর্ণ: ক্রমসম্পর্ক এবং তুল্যতা সম্পর্ক।
\end{intro}

\begin{defn}[প্রতিবিম্ব ধর্ম]
একটি সম্পর্ক $R \subseteq A^2$ \emph{প্রতিবিম্ব ধর্মবিশিষ্ট} যদি এবং কেবল যদি প্রত্যেক $x \in A$-এর জন্য $Rxx$ হয়।
\end{defn}

\begin{defn}[পরিযায়িতা]
একটি সম্পর্ক $R \subseteq A^2$ \emph{পরিযায়ী} যদি এবং কেবল যদি যখনই $Rxy$ এবং $Ryz$ হয়, তখনই $Rxz$-ও হয়।
\end{defn}

\begin{defn}[প্রতিসাম্য]
একটি সম্পর্ক~$R \subseteq A^2$ \emph{প্রতিসম} যদি এবং কেবল যদি যখনই $Rxy$ হয়, তখনই~$Ryx$-ও হয়।
\end{defn}

\begin{defn}[বিপ্রতিসাম্য]
একটি সম্পর্ক~$R \subseteq A^2$ \emph{বিপ্রতিসম} যদি এবং কেবল যদি $Rxy$ ও $Ryx$ উভয়ই হলে $x=y$ হয় (অন্যভাবে বললে: $x\neq y$ হলে হয় $\lnot Rxy$, নয়তো $\lnot Ryx$)।
\end{defn}

\begin{explain}
প্রতিসম সম্পর্কে $Rxy$ ও $Ryx$ সর্বদা একসঙ্গে সত্য হয়, অথবা কোনোটিই সত্য হয় না। বিপ্রতিসম সম্পর্কে $Rxy$ ও $Ryx$ একসঙ্গে সত্য হওয়ার একমাত্র সম্ভাবনা হল $x = y$। লক্ষ করো, এর অর্থ এই নয় যে $x = y$ হলে $Rxy$ ও $Ryx$ \emph{অবশ্যই} সত্য হবে; কেবল এটুকুই যে তা নিষিদ্ধ নয়। তাই বিপ্রতিসম সম্পর্ক প্রতিবিম্ব ধর্মবিশিষ্ট হতে পারে, কিন্তু প্রত্যেক বিপ্রতিসম সম্পর্ক প্রতিবিম্ব ধর্মবিশিষ্ট নয়। আরও লক্ষ করো, বিপ্রতিসম হওয়া এবং কেবল প্রতিসম না হওয়া ভিন্ন শর্ত। বস্তুত, একটি সম্পর্ক একই সঙ্গে প্রতিসম ও বিপ্রতিসম হতে পারে (যেমন, অভিন্নতার সম্পর্ক)।
\end{explain}

\begin{defn}[সংযুক্ততা]
একটি সম্পর্ক $R \subseteq A^2$ \emph{সংযুক্ত} যদি সকল $x,y\in A$-এর জন্য, $x \neq y$ হলে হয় $Rxy$, নয়তো~$Ryx$ হয়।
\end{defn}

\begin{prob}
এমন সম্পর্কের উদাহরণ দাও যা (a) প্রতিবিম্ব ধর্মবিশিষ্ট ও প্রতিসম কিন্তু পরিযায়ী নয়, (b) প্রতিবিম্ব ধর্মবিশিষ্ট ও বিপ্রতিসম, (c) বিপ্রতিসম ও পরিযায়ী কিন্তু প্রতিবিম্ব ধর্মবিশিষ্ট নয়, এবং (d) প্রতিবিম্ব ধর্মবিশিষ্ট, প্রতিসম ও পরিযায়ী। সংখ্যা বা সেটের উপর সম্পর্ক ব্যবহার কোরো না।
\end{prob}

\begin{defn}[আত্মসম্পর্কহীনতা]
একটি সম্পর্ক $R \subseteq A^2$-কে \emph{আত্মসম্পর্কহীন} বলা হয় যদি সকল $x \in A$-এর জন্য $Rxx$ সত্য না হয়।
\end{defn}

\begin{defn}[একমুখিতা]
একটি সম্পর্ক $R \subseteq A^2$-কে \emph{একমুখী} বলা হয় যদি কোনো $x,y\in A$ যুগলের জন্যই $Rxy$ ও~$Ryx$ একসঙ্গে সত্য না হয়।
\end{defn}

লক্ষ করো, $A \neq \emptyset$ হলে $A$-এর উপর কোনো আত্মসম্পর্কহীন সম্পর্ক প্রতিবিম্ব ধর্মবিশিষ্ট হয় না, এবং $A$-এর উপর প্রত্যেক একমুখী সম্পর্ক বিপ্রতিসমও হয়। কিন্তু এমন $R \subseteq A^2$ আছে যা প্রতিবিম্ব ধর্মবিশিষ্টও নয়, আত্মসম্পর্কহীনও নয়; আবার এমন বিপ্রতিসম সম্পর্ক আছে যা একমুখী নয়।

\end{document}
